\section{Электричество}
\subsection{Закон Кулона}

%%%%%%%%%%%%%%%%%%%%%%%%%%%%%%%%%%%%%%%%%%%%%%%%%%%%
%%%%%%%%%%%%%%%%%%%%%%%%%%%%%%%%%%%%%%%%%%%%%%%%%%%%
 \z{Заряженный шарик приводят в соприкосновение с точно таким же незаряженным шариком. Находясь на расстоянии $r = 15\,\text{см}$, шарики отталкиваются с силой $F = 1\,\text{мН}$. Каков первоначальный заряд шарика?}
%
{$q = \sqrt{\frac{4Fr^2}{k}} = 10^{-7}\,\text{Кл}$}

%%%%%%%%%%%%%%%%%%%%%%%%%%%%%%%%%%%%%%%%%%%%%%%%%%%%
 \z{Две отрицательно заряженные пылинки находятся в воздухе на расстоянии $r = 2\,\text{мм}$ друг от друга и отталкиваются с силой $F = 9 \cdot 10^{-5}\,\text{Н}$. Считая заряды пылинок одинаковыми, найти число избыточных электронов на каждой пылинке.}
%
{$N = \sqrt{\frac{Fr^2}{ke^2}} = 1.25 \cdot 10^9\,\text{шт}$}

%%%%%%%%%%%%%%%%%%%%%%%%%%%%%%%%%%%%%%%%%%%%%%%%%%%%
 \z[{\begin{wrapfigure}[3]{r}{0.3\textwidth}%
		\vspace{-15pt}%
		\begin{flushright}%
 			\includegraphics[width=0.3\textwidth]{3.1.3}%
		\end{flushright}%
	\end{wrapfigure}%
	}]%
{Четыре маленьких одинаковых шарика, связанных нерастяжимыми нитями одинаковой длины, заряжены зарядами $q$, $q$, $q$ и $2q$. Сила натяжения нити, связывающей первый и второй шарики, равна $T$. Найти силу натяжения нити, связывающей второй и третий шарики.}
%
{$T_{23} = \frac{71}{53}T$}

%%%%%%%%%%%%%%%%%%%%%%%%%%%%%%%%%%%%%%%%%%%%%%%%%%%%
 \z[{\begin{wrapfigure}[5]{r}{0.3\textwidth}%
		\vspace{-15pt}%
		\begin{flushright}%
 			\includegraphics[width=0.3\textwidth]{3.1.4}%
		\end{flushright}%
	\end{wrapfigure}%
	}]%
 {Две маленьких бусинки массой $m$ заряжены зарядами $Q$ и $Q$. Бусинки надеты на спицы, которые расположены в вертикальной плоскости симметрично по отношению к вертикали, и угол между которыми равен $90^\circ$ (см. рисунок). Каково расстояние между бусинками в положении равновесия?}
%
{$d = \sqrt{\frac{kQ^2}{mg}}$}

%%%%%%%%%%%%%%%%%%%%%%%%%%%%%%%%%%%%%%%%%%%%%%%%%%%%
 \z{Два электрона находятся в точках, определяемых радиус-векторами $(2;2)$ и $(3;-2)$. Вычислить ускорение, сообщаемое одним электроном другому.}
%
{$a = 15\,\text{м/с}^2$}

%%%%%%%%%%%%%%%%%%%%%%%%%%%%%%%%%%%%%%%%%%%%%%%%%%%%
 \z[{\begin{wrapfigure}[7]{r}{0.3\textwidth}%
		\vspace{-15pt}%
		\begin{flushright}%
 			\includegraphics[width=0.3\textwidth]{3.1.6}%
		\end{flushright}%
	\end{wrapfigure}%
	}]%
 {Три точечных заряда $q = 10\,\text{мкКл}$, $2q$ и $-3q$ расположены на окружности радиуса $R = 30\,\text{см}$. Найти модуль и направление силы, действующей на заряд $q$ со стороны двух других.}
%
{$F=\frac{kq^2}{(2R)^2} \sqrt{13} = 18\,\text{Н}$}

%%%%%%%%%%%%%%%%%%%%%%%%%%%%%%%%%%%%%%%%%%%%%%%%%%%%
 \z{На гладкую замкнутую непроводящую нить длиной $l = 60\,\text{см}$ нанизаны три бусинки с зарядами $q_1 = 20\,\text{мкКл}$, $q_2 = q_1/4$, $q_3 = q_1/9$. Система находится в равновесии. Найти силу натяжения нити $T$.}
%
{$T = \frac{kq^2}{l^2} = 10\,\text{Н}$}

